\documentclass[11pt]{article}
\usepackage[margin=1in]{geometry}
\usepackage{amsmath,amssymb,amsthm}
\usepackage{booktabs}
\usepackage[hidelinks]{hyperref}

\newtheorem{proposition}{Proposition}
\newtheorem{lemma}{Lemma}
\newtheorem{corollary}{Corollary}
\theoremstyle{remark}
\newtheorem{remark}{Remark}

\allowdisplaybreaks

\title{The Beveridge Response to a Destruction Shock\\ under Infinitely Elastic Vacancy Creation}
\author{Notes for review (not for the draft)}
\date{October 6, 2026}

\begin{document}
\maketitle

\noindent Numbers are produced by \texttt{Programs baseline/gs\_free\_entry\_beveridge.jl}
(output \texttt{gs\_free\_entry\_beveridge.txt}, block ``Exposition tables''). The program
checks every closed form below by assertion against a simulated impulse response of the
same linear system.

\section{Setting}

We work with the model under the conditions of the draft's independence proposition
(\texttt{prop:independence}): $\sigma=0$, $\zeta=0$ and $p_0=0$. The labor-market block
then closes on its own. It is the Gabrovski--Silva economy with a time-varying destruction
rate $\delta_t$. In this block:
\begin{itemize}
	\item the discount factor is $\beta$, and $r\equiv1/\beta-1$ is the discount rate;
	\item recruiter compensation $y\equiv z/\mu$ is constant;
	\item exit is exogenous, so a business survives from $t$ to $t+1$ with probability
	$1-\delta_t$;
	\item every vacated position is reposted at no cost.
\end{itemize}
For exposition we set the per-match cost $\kappa=0$, so all recruiting costs are sunk
($x_v=1$). Remark~\ref{rem:kappa} discusses $\kappa>0$.

\paragraph{Timing}
We use the draft's convention. $\delta_t$ is known at the start of $t$ and destroys
businesses at the end of $t$. Hence $u_t$ and the inherited vacancy stock are fixed at $t$.
Tightness $\theta_t$ and entry $e_t$ respond at $t$. The destroyed jobs and positions
appear in $u_{t+1}$ and $v_{t+1}$.

\paragraph{Matching}
$q(\theta)=A\theta^{-\eta_L}$ and $f(\theta)=\theta q(\theta)=A\theta^{1-\eta_L}$, with
$\eta_L\in(0,1)$.

\section{The nonlinear system}

\subsection{The vacancy value and its flow value under free entry}

New vacancies require a sunk cost drawn from $[0,x_m]$. With infinitely elastic creation,
the value of a vacancy is pinned whenever entry is positive:
\begin{equation}\label{eq:n_fe}
	Q_t=x_m\quad\text{for all }t.
\end{equation}
The flow value of holding a vacancy for one period (eq:Kdef in the draft) is
$K_t=Q_t-\beta(1-\delta_t)Q_{t+1}$. The firm holds an asset worth $Q_t$, and it is worth
$Q_{t+1}$ next period only if the business survives destruction at the end of $t$.
Substituting \eqref{eq:n_fe},
\begin{equation}\label{eq:n_K}
	K_t=x_m\bigl[1-\beta(1-\delta_t)\bigr]=x_m(1-\beta+\beta\delta_t).
\end{equation}
\textbf{$Q$ is a constant. $K_t$ is a known function of the exogenous rate $\delta_t$,
not an endogenous variable.} When $\delta$ is constant, as in Gabrovski--Silva, $K$ is
constant too. When $\delta_t$ moves, $K_t$ moves with it, because a vacancy is an asset
and $\delta_t$ is its depreciation rate. Holding $K$ fixed instead would give the
flow-cost DMP model (Remark~\ref{rem:flow}). We substitute \eqref{eq:n_K} wherever $K$
appears, so $K$ is not carried as a variable.

\subsection{The remaining equations}

\begin{align}
	&\frac{x_m(1-\beta+\beta\delta_t)}{q(\theta_t)}
	=\beta(1-\delta_t)\,\mathbb{E}_t\Bigl[(1-\phi)\bigl(y-b-x_m(1-\beta+\beta\delta_{t+1})\bigr)
	\notag\\
	&\hspace{9em}+(1-s)\frac{x_m(1-\beta+\beta\delta_{t+1})}{q(\theta_{t+1})}
	-\phi\,\theta_{t+1}\,x_m(1-\beta+\beta\delta_{t+1})\Bigr], \label{eq:n_jcc}\\
	&u_{t+1}=u_t-(1-\delta_t)f(\theta_t)\,u_t+\bigl[\delta_t+s(1-\delta_t)\bigr](1-u_t),
	\label{eq:n_u}\\
	&v_{t+1}=(1-\delta_t)\bigl[(1-q(\theta_t))\,v_t+s(1-u_t)\bigr]+e_{t+1}, \label{eq:n_v}\\
	&\theta_t=v_t/u_t. \label{eq:n_theta}
\end{align}

\paragraph{\eqref{eq:n_jcc} Job creation}
\begin{itemize}
	\item \emph{The left side is the expected cost of a hire.} The vacancy's flow value
	$x_m(1-\beta+\beta\delta_t)$ is paid for an expected $1/q(\theta_t)$ periods.
	\item \emph{The right side is the value of a match.} The match first produces at $t+1$,
	and only if the business survives.
	\item \emph{The bracket is the firm's surplus after substituting the Nash wage.} It has
	three terms:
	\begin{itemize}
		\item The first is the firm's share of output, net of the worker's outside option and
		of the vacancy's flow value, which a matched firm forgoes.
		\item The second is the hiring cost saved next period because the match continues.
		\item The third is the part of that saving the worker captures through the wage. The
		worker's outside option includes finding another job, and $\phi f K/q=\phi\theta K$
		because $f(\theta)/q(\theta)=\theta$.
	\end{itemize}
\end{itemize}

\paragraph{\eqref{eq:n_u} Unemployment}
Next period's unemployed are:
\begin{itemize}
	\item today's unemployed, less those hired at a business that survives;
	\item plus the employed who lose their job, to destruction ($\delta_t$) or to separation
	at a surviving business ($s(1-\delta_t)$).
\end{itemize}

\paragraph{\eqref{eq:n_v} Vacancies}
Next period's vacancies are:
\begin{itemize}
	\item today's unfilled vacancies that survive destruction;
	\item positions vacated by separations at surviving businesses and reposted at no cost;
	\item new entry.
\end{itemize}

\paragraph{How free entry orders the system}
Equation \eqref{eq:n_jcc} contains only $\theta$ and the exogenous $\delta$. Given the
path of $\delta$, it determines the path of $\theta$ on its own. Then:
\begin{itemize}
	\item \eqref{eq:n_u} gives $u$;
	\item \eqref{eq:n_theta} gives $v=\theta u$;
	\item \eqref{eq:n_v} only determines entry $e_{t+1}$, as whatever delivers
	$v_{t+1}=\theta_{t+1}u_{t+1}$ given the inherited stock.
\end{itemize}

\section{Steady state}

Write
\[
	\bar K\equiv x_m(1-\beta+\beta\bar\delta)=x_m\,\frac{r+\bar\delta}{1+r},
	\qquad
	\bar\tau\equiv\bar\delta+\bar s(1-\bar\delta).
\]
$\bar K$ is the steady-state flow value of a vacancy. It is a user cost: the asset price
times the discount rate plus the destruction rate. The steady-state versions of
\eqref{eq:n_jcc} and \eqref{eq:n_u} are
\begin{align}
	\frac{\bar K}{\bar q}
	&=\beta(1-\bar\delta)\Bigl[(1-\phi)(y-b-\bar K)+(1-\bar s)\frac{\bar K}{\bar q}
	-\phi\bar\theta\bar K\Bigr], \label{eq:ss_jcc}\\
	\bar u&=\frac{\bar\tau}{\bar\tau+(1-\bar\delta)\bar f},
	\qquad\text{equivalently}\qquad
	\frac{1-\bar u}{\bar u}=\frac{(1-\bar\delta)\bar f}{\bar\tau}. \label{eq:ss_u}
\end{align}

\section{Log-linear system}

Hats denote log deviations: $\hat x_t\equiv\ln x_t-\ln\bar x$. For the destruction rate we
use the level deviation $d_t\equiv\delta_t-\bar\delta$.

\subsection{Job creation}

Take logs of both sides of \eqref{eq:n_jcc}.

\paragraph{Left side}
\[
	\ln\frac{x_m(1-\beta+\beta\delta_t)}{q(\theta_t)}
	=\ln x_m+\ln(1-\beta+\beta\delta_t)-\ln A+\eta_L\ln\theta_t .
\]
Its first-order deviation is
\[
	\frac{\beta}{1-\beta+\beta\bar\delta}\,d_t+\eta_L\hat\theta_t
	=\frac{d_t}{r+\bar\delta}+\eta_L\hat\theta_t,
\]
using $\beta/(1-\beta+\beta\bar\delta)=1/(r+\bar\delta)$.

\paragraph{Right side}
The log of $\beta(1-\delta_t)$ deviates by $-d_t/(1-\bar\delta)$. Write $\Sigma_{t+1}$ for
the bracket in \eqref{eq:n_jcc}. Its steady-state value $\bar\Sigma$ satisfies
$\bar K/\bar q=\beta(1-\bar\delta)\bar\Sigma$ by \eqref{eq:ss_jcc}. The deviation of
$\Sigma_{t+1}$ is the sum of three terms:
\begin{align*}
	d\Bigl[(1-\phi)\bigl(y-b-x_m(1-\beta+\beta\delta_{t+1})\bigr)\Bigr]
	&=-(1-\phi)\,\beta x_m\,d_{t+1}
	=-(1-\phi)\,\bar K\,\frac{d_{t+1}}{r+\bar\delta},\\
	d\Bigl[(1-s)\frac{x_m(1-\beta+\beta\delta_{t+1})}{q(\theta_{t+1})}\Bigr]
	&=(1-\bar s)\frac{\bar K}{\bar q}\Bigl(\frac{d_{t+1}}{r+\bar\delta}+\eta_L\hat\theta_{t+1}\Bigr),\\
	d\Bigl[-\phi\,\theta_{t+1}x_m(1-\beta+\beta\delta_{t+1})\Bigr]
	&=-\phi\,\bar\theta\bar K\Bigl(\frac{d_{t+1}}{r+\bar\delta}+\hat\theta_{t+1}\Bigr).
\end{align*}
Divide the sum by $\bar\Sigma=\bar K/[\bar q\,\beta(1-\bar\delta)]$, and use
$\bar\theta\bar q=\bar f$:
\[
	\hat\Sigma_{t+1}=\beta(1-\bar\delta)\Bigl[-(1-\phi)\bar q\,\frac{d_{t+1}}{r+\bar\delta}
	+(1-\bar s)\Bigl(\frac{d_{t+1}}{r+\bar\delta}+\eta_L\hat\theta_{t+1}\Bigr)
	-\phi\bar f\Bigl(\frac{d_{t+1}}{r+\bar\delta}+\hat\theta_{t+1}\Bigr)\Bigr].
\]

\paragraph{Equating the two sides}
Set the left-side deviation equal to $-d_t/(1-\bar\delta)+\mathbb{E}_t\hat\Sigma_{t+1}$.
Write $B\equiv\beta(1-\bar\delta)$, move the $d_t$ term to the right, and collect terms:
\begin{equation}\label{eq:l_jcc}
	\boxed{\;
	\eta_L\hat\theta_t
	=-\Bigl[\frac{1}{r+\bar\delta}+\frac{1}{1-\bar\delta}\Bigr]d_t
	+B\,\mathbb{E}_t\Bigl\{\bigl[1-\bar s-\phi\bar f-(1-\phi)\bar q\bigr]\frac{d_{t+1}}{r+\bar\delta}
	+\bigl[(1-\bar s)\eta_L-\phi\bar f\bigr]\hat\theta_{t+1}\Bigr\}\;}
\end{equation}
Only $\hat\theta$ and the shock appear.

\paragraph{Interpretation of each weight}
\begin{itemize}
	\item \emph{$\eta_L$ on $\hat\theta_t$.} A tighter market lengthens the time to fill, so
	a hire costs more.
	\item \emph{$1/(r+\bar\delta)$ on $d_t$: the rental effect.} A one-point rise in
	destruction raises the vacancy's flow value $x_m(1-\beta+\beta\delta_t)$ by $\beta x_m$,
	which is $1/(r+\bar\delta)$ in proportion. The longer-lived the vacancy, the larger the
	proportional rise. This raises the cost of a hire.
	\item \emph{$1/(1-\bar\delta)$ on $d_t$: the survival effect.} A hire pays off only if
	the business survives. This lowers the value of a hire. It is about 1.
	\item \emph{On $d_{t+1}$, which is nonzero only if the shock persists:}
	\begin{itemize}
		\item $1-\bar s-\phi\bar f$: a higher flow value next period raises the value of an
		ongoing match, which saves next period's hiring cost, net of the worker's share;
		\item $-(1-\phi)\bar q$: a matched firm forgoes that flow value, so the firm's surplus
		falls.
	\end{itemize}
	\item \emph{On $\hat\theta_{t+1}$:}
	\begin{itemize}
		\item $(1-\bar s)\eta_L$: a tighter market next period raises the hiring cost an
		ongoing match saves;
		\item $-\phi\bar f$: it also raises the worker's outside option and hence the wage.
	\end{itemize}
\end{itemize}

\subsection{Unemployment}

Write the right side of \eqref{eq:n_u} as
$U(u_t,\delta_t,\theta_t)=u_t-(1-\delta_t)A\theta_t^{1-\eta_L}u_t+[\delta_t+s(1-\delta_t)](1-u_t)$.
Its partial derivatives at the steady state are
\begin{align*}
	\frac{\partial U}{\partial u_t}&=1-(1-\bar\delta)\bar f-\bigl[\bar\delta+\bar s(1-\bar\delta)\bigr]
	=1-(1-\bar\delta)\bar f-\bar\tau,\\
	\frac{\partial U}{\partial\delta_t}&=\bar f\,\bar u+(1-\bar s)(1-\bar u),\\
	\frac{\partial U}{\partial\theta_t}&=-(1-\bar\delta)\,\bar u\,(1-\eta_L)\frac{\bar f}{\bar\theta}.
\end{align*}
Hence
\[
	u_{t+1}-\bar u
	=\bigl[1-(1-\bar\delta)\bar f-\bar\tau\bigr](u_t-\bar u)
	+\bigl[\bar f\bar u+(1-\bar s)(1-\bar u)\bigr]d_t
	-(1-\bar\delta)\bar u(1-\eta_L)\bar f\,\hat\theta_t .
\]
Divide by $\bar u$. The coefficient on $d_t$ becomes, by \eqref{eq:ss_u},
\[
	\frac{\bar f\bar u+(1-\bar s)(1-\bar u)}{\bar u}
	=\bar f+(1-\bar s)\frac{(1-\bar\delta)\bar f}{\bar\tau}
	=\frac{\bar f}{\bar\tau}\bigl[\bar\tau+(1-\bar s)(1-\bar\delta)\bigr]
	=\frac{\bar f}{\bar\tau},
\]
because $(1-\bar s)(1-\bar\delta)=1-\bar\tau$. Define
\[
	\lambda\equiv1-(1-\bar\delta)\bar f-\bar\tau,
	\qquad
	\omega\equiv(1-\bar\delta)\bar f(1-\eta_L).
\]
Then
\begin{equation}\label{eq:l_u}
	\boxed{\;\hat u_{t+1}=\lambda\,\hat u_t+\frac{\bar f}{\bar\tau}\,d_t-\omega\,\hat\theta_t\;}
\end{equation}

\paragraph{Interpretation of each weight}
\begin{itemize}
	\item \emph{$\lambda$ is persistence.} An extra unemployed worker is still unemployed
	next month unless hired, with probability $1-(1-\bar\delta)\bar f$. The $\bar\tau$ term
	reflects that one fewer employed worker means one fewer separation.
	\item \emph{$\bar f/\bar\tau$ is the destruction weight.} A one-point rise in $\delta_t$
	adds two groups to unemployment: the employed who lose their job, $(1-\bar s)(1-\bar u)$,
	and the new matches lost before producing, $\bar f\bar u$. Relative to the stock $\bar u$
	this is $\bar f/\bar\tau$. It is large because unemployment is a small stock, of order
	$\bar\tau$. It depends on $\bar f$ and $\bar\tau$, not on how $\bar\tau$ splits into
	$\delta$ and $s$.
	\item \emph{$\omega$ is the hiring channel.} Lower tightness at $t$ means fewer hires at
	$t$, who are still unemployed at $t+1$.
\end{itemize}

\subsection{Vacancies}

Write the right side of \eqref{eq:n_v} as
$V(v_t,u_t,\delta_t,\theta_t,e_{t+1})
=(1-\delta_t)[(1-A\theta_t^{-\eta_L})v_t+s(1-u_t)]+e_{t+1}$. Its partial derivatives at the
steady state are
\begin{align*}
	\frac{\partial V}{\partial\delta_t}&=-\bigl[(1-\bar q)\bar v+\bar s(1-\bar u)\bigr]\equiv-\bar X,
	&\frac{\partial V}{\partial v_t}&=(1-\bar\delta)(1-\bar q),\\
	\frac{\partial V}{\partial\theta_t}&=(1-\bar\delta)\,\eta_L\frac{\bar q}{\bar\theta}\,\bar v,
	&\frac{\partial V}{\partial u_t}&=-(1-\bar\delta)\bar s,
	\qquad\frac{\partial V}{\partial e_{t+1}}=1 .
\end{align*}
Hence
\[
	v_{t+1}-\bar v=-\bar X d_t+(1-\bar\delta)(1-\bar q)(v_t-\bar v)
	+(1-\bar\delta)\eta_L\bar q\,\bar v\,\hat\theta_t
	-(1-\bar\delta)\bar s\,(u_t-\bar u)+(e_{t+1}-\bar e).
\]
Divide by $\bar v$, and use $\bar u/\bar v=1/\bar\theta$:
\begin{equation}\label{eq:l_v}
	\boxed{\;
	\hat v_{t+1}=-\frac{\bar X}{\bar v}\,d_t
	+(1-\bar\delta)(1-\bar q)\,\hat v_t
	+(1-\bar\delta)\eta_L\bar q\,\hat\theta_t
	-(1-\bar\delta)\frac{\bar s}{\bar\theta}\,\hat u_t
	+\frac{\bar e}{\bar v}\,\hat e_{t+1}\;}
\end{equation}

\paragraph{Interpretation of each weight}
\begin{itemize}
	\item \emph{$\bar X/\bar v$.} This is the share of next period's stock inherited from
	this period. Destruction removes it one for one.
	\item \emph{$(1-\bar\delta)(1-\bar q)$.} This is the share of vacancies that are neither
	filled nor destroyed.
	\item \emph{$(1-\bar\delta)\eta_L\bar q$.} A tighter market fills fewer vacancies, so
	more carry over. After a destruction shock $\hat\theta<0$, so vacancies fill faster.
	\item \emph{$(1-\bar\delta)\bar s/\bar\theta$.} Reposts come from the employed, who are
	fewer when $u$ is higher.
	\item \emph{$\bar e/\bar v$.} This is entry's share of the stock.
\end{itemize}
At free entry, \eqref{eq:l_v} determines $\hat e_{t+1}$, not $\hat v_{t+1}$.

\subsection{Tightness}

Taking logs of \eqref{eq:n_theta},
\begin{equation}\label{eq:l_theta}
	\boxed{\;\hat v_t=\hat\theta_t+\hat u_t\;}
\end{equation}

\subsection{Weights at the calibration}

Our calibration holds $\bar f$, $\bar q$ and $\bar\tau$ fixed and recalibrates $\bar s$ and
$\phi$ at each $\bar\delta$.

\begin{table}[h]
\centering
\begin{tabular}{lcc}
\toprule
 & D1 ($\bar\delta=3.2\%$/yr) & CK-like ($\bar\delta=\bar\tau$) \\
\midrule
$1/(r+\bar\delta)$ & 167.2 & 29.2 \\
$\bar f/\bar\tau$ & 13.26 & 13.65 \\
$\lambda$ & 0.559 & 0.559 \\
$\omega$ & 0.164 & 0.164 \\
$\bar X/\bar v$ & 0.930 & 0.174 \\
$\bar e/\bar v$ & 0.073 & 0.831 \\
\bottomrule
\end{tabular}
\caption{Weights of the log-linear system.}
\label{tab:weights}
\end{table}

The unemployment equation is essentially the same in both economies. Among the weights
that matter for $u$ and $v$ at free entry, only $1/(r+\bar\delta)$ differs much.

\section{Tightness}

\begin{lemma}\label{lem:theta}
Let $d_t=\rho^t d_0$ with $\rho\in[0,1)$, and suppose
\[
	\Phi\equiv B\Bigl[(1-\bar s)-\frac{\phi\bar f}{\eta_L}\Bigr]>-1 .
\]
Then \eqref{eq:l_jcc} has a unique bounded solution, $\hat\theta_t=a\,d_t$, with
\begin{equation}\label{eq:a}
	a=-\,\frac{\dfrac{1}{1-\bar\delta}
	+\dfrac{1-B\rho\bigl[1-\bar s-\phi\bar f-(1-\phi)\bar q\bigr]}{r+\bar\delta}}
	{\eta_L-B\rho\bigl[(1-\bar s)\eta_L-\phi\bar f\bigr]}<0 .
\end{equation}
\end{lemma}

\begin{proof}
\emph{Uniqueness.} Divide \eqref{eq:l_jcc} by $\eta_L$:
\[
	\hat\theta_t=\Phi\,\mathbb{E}_t\hat\theta_{t+1}
	-\frac{1}{\eta_L}\Bigl[\frac{1}{r+\bar\delta}+\frac{1}{1-\bar\delta}\Bigr]d_t
	+\frac{B\bigl[1-\bar s-\phi\bar f-(1-\phi)\bar q\bigr]}{\eta_L(r+\bar\delta)}\,
	\mathbb{E}_td_{t+1}.
\]
Here $\Phi<B(1-\bar s)<1$, and $\Phi>-1$ by assumption. So the unique bounded solution is
obtained by solving forward. With $\mathbb{E}_td_{t+k}=\rho^kd_t$, every term of the
forward sum is proportional to $d_t$, so $\hat\theta_t=a\,d_t$ for some constant $a$.

\emph{The coefficient.} Substitute $\hat\theta_t=a\rho^td_0$,
$\mathbb{E}_t\hat\theta_{t+1}=a\rho^{t+1}d_0$, $d_t=\rho^td_0$ and
$\mathbb{E}_td_{t+1}=\rho^{t+1}d_0$ into \eqref{eq:l_jcc}:
\[
	\eta_L a\rho^td_0
	=-\Bigl[\frac{1}{r+\bar\delta}+\frac{1}{1-\bar\delta}\Bigr]\rho^td_0
	+B\Bigl\{\bigl[1-\bar s-\phi\bar f-(1-\phi)\bar q\bigr]\frac{\rho^{t+1}d_0}{r+\bar\delta}
	+\bigl[(1-\bar s)\eta_L-\phi\bar f\bigr]a\rho^{t+1}d_0\Bigr\}.
\]
Divide by $\rho^td_0$:
\[
	\eta_L a=-\frac{1}{r+\bar\delta}-\frac{1}{1-\bar\delta}
	+\frac{B\rho\bigl[1-\bar s-\phi\bar f-(1-\phi)\bar q\bigr]}{r+\bar\delta}
	+B\rho\bigl[(1-\bar s)\eta_L-\phi\bar f\bigr]a .
\]
Collect the terms in $a$ on the left:
\[
	a\Bigl\{\eta_L-B\rho\bigl[(1-\bar s)\eta_L-\phi\bar f\bigr]\Bigr\}
	=-\frac{1}{1-\bar\delta}
	-\frac{1-B\rho\bigl[1-\bar s-\phi\bar f-(1-\phi)\bar q\bigr]}{r+\bar\delta},
\]
which is \eqref{eq:a}.

\emph{The sign.} $B\rho<1$ and $1-\bar s-\phi\bar f-(1-\phi)\bar q<1$, so the numerator of
the second fraction is positive. The denominator satisfies
$\eta_L-B\rho[(1-\bar s)\eta_L-\phi\bar f]\ge\eta_L[1-B\rho(1-\bar s)]>0$. Hence $a<0$.
\end{proof}

\begin{remark}[Why the endogenous states do not enter]\label{rem:recursive}
The economy has endogenous states, $u_t$ and the predetermined vacancies $v_{pre,t}$. A
general linear solution for tightness would take the form
$\hat\theta_t=a\,d_t+b_u\hat u_t+b_v\hat v_{pre,t}$. The proof above nevertheless delivers
$b_u=b_v=0$, for the following reasons.
\begin{enumerate}
	\item Equation \eqref{eq:l_jcc} contains neither $u$ nor $v_{pre}$. Free entry pins $Q$,
	and $K_t=x_m(1-\beta+\beta\delta_t)$ depends only on the exogenous $\delta_t$.
	\item Iterating \eqref{eq:l_jcc} forward $T$ periods gives
	\[
		\hat\theta_t=\sum_{k=0}^{T-1}\Phi^k\,\mathbb{E}_t\bigl(c_0d_{t+k}+c_1d_{t+k+1}\bigr)
		+\Phi^T\,\mathbb{E}_t\hat\theta_{t+T},
	\]
	where $c_0$ and $c_1$ are the coefficients in the proof. For a bounded solution with
	$\lvert\Phi\rvert<1$, the last term vanishes as $T\to\infty$. So $\hat\theta_t$ depends
	only on current and expected future $d$, which gives $b_u=b_v=0$.
	\item The states are then solved backward. Iterating \eqref{eq:l_u} gives $\hat u$
	(Proposition~\ref{prop:fe}, Part 2), and \eqref{eq:l_theta} gives
	$\hat v=\hat\theta+\hat u$.
	\item The predetermined stock
	$v_{pre,t+1}=(1-\delta_t)[(1-q(\theta_t))\theta_tu_t+s(1-u_t)]$ depends only on
	$(u_t,\theta_t,\delta_t)$ and has no own lag. It feeds into nothing except entry,
	$e_{t+1}=v_{t+1}-v_{pre,t+1}$.
\end{enumerate}
\emph{Root count.} The system has one unstable root, $1/\Phi$, for one jump variable, and
one stable root, $\lambda$, for one dynamic state, $u$. So the Blanchard--Kahn condition
holds.

\emph{When the structure breaks.} With finite $\xi$, entry $e=G(Q)$ links $v_{pre}$ to $Q$,
and $Q$ enters the job creation condition. Tightness then depends on the states, the guess
must include state coefficients, and the closed form is lost. The same happens when the
conditions of \texttt{prop:independence} fail, through $N$, $C$ or $\chi^c$.
\end{remark}

\paragraph{Reading \eqref{eq:a}}
\begin{itemize}
	\item $1/(1-\bar\delta)$ is the survival effect.
	\item $1/(r+\bar\delta)$ is the rental effect.
	\item Persistence adds two offsetting adjustments, both because the flow value stays
	high next period:
	\begin{itemize}
		\item it raises the value of an ongoing match, through $1-\bar s-\phi\bar f$;
		\item it erodes next period's surplus, through $-(1-\phi)\bar q$.
	\end{itemize}
	\item The denominator accounts for tightness being low next period as well.
	\item At $\rho=0$,
	$a=-\bigl[1/(r+\bar\delta)+1/(1-\bar\delta)\bigr]/\eta_L$.
\end{itemize}

\section{Proposition}

\begin{proposition}[Destruction shock under infinitely elastic vacancy creation]
\label{prop:fe}
Assume the setting of Section~1, with $\xi\to\infty$, the condition $\Phi>-1$ of
Lemma~\ref{lem:theta}, and $\lambda\in[0,1)$. Let $\delta_t=\bar\delta+\rho^td_0$ with
$d_0>0$ and $\rho\in[0,1)$, and let $a<0$ be given by \eqref{eq:a}. To first order:
\begin{enumerate}
	\item \emph{Impact.} $\hat u_0=0$ and $\hat v_0=\hat\theta_0=a\,d_0<0$.
	\item \emph{Unemployment.} For $t\ge1$,
	\[
		\hat u_t=\Bigl(\frac{\bar f}{\bar\tau}+\omega\lvert a\rvert\Bigr)d_0\,g_t>0,
		\qquad g_t\equiv\sum_{j=0}^{t-1}\lambda^{t-1-j}\rho^j .
	\]
	\item \emph{Vacancies at $h=1$.} $\hat v_1<0$ if and only if
	\begin{equation}\label{eq:H1}
		\lvert a\rvert(\rho-\omega)>\frac{\bar f}{\bar\tau}.
	\end{equation}
	\item \emph{Vacancies at later horizons.} If $\hat v_t\ge0$ for some $t\ge1$, then
	$\hat v_{t'}>0$ for all $t'>t$. Moreover, $\hat v_t<0$ for every $t\ge0$ if and only if
	\[
		\rho>\lambda\quad\text{and}\quad\lvert a\rvert(\rho-\lambda-\omega)\ge\frac{\bar f}{\bar\tau}.
	\]
	\item \emph{The whole path.} The slope of a regression of $\tilde v_t\equiv v_t-\bar v$ on
	$\tilde u_t\equiv u_t-\bar u$ along the impulse response ($t\ge0$) is
	\begin{equation}\label{eq:slope}
		\frac{\sum_t\tilde v_t\tilde u_t}{\sum_t\tilde u_t^2}
		=\bar\theta\left[1-\frac{\lvert a\rvert\,\rho(1-\lambda^2)}
		{\bigl(\bar f/\bar\tau+\omega\lvert a\rvert\bigr)(1+\rho\lambda)}\right].
	\end{equation}
	It is negative if and only if
	\begin{equation}\label{eq:HP}
		\lvert a\rvert\bigl[\rho(1-\lambda^2)-\omega(1+\rho\lambda)\bigr]
		>\frac{\bar f}{\bar\tau}(1+\rho\lambda).
	\end{equation}
\end{enumerate}
\end{proposition}

\begin{corollary}[iid shock]\label{cor:iid}
If $\rho=0$, then $\hat\theta_t=0$ and $\hat v_t=\hat u_t>0$ for all $t\ge1$. Unemployment
and vacancies move together along the ray $v=\bar\theta u$ at every horizon after impact,
and the slope \eqref{eq:slope} equals $\bar\theta>0$.
\end{corollary}

\begin{proof}[Proof of Proposition~\ref{prop:fe}]
\emph{Part 1.} $u_0$ is predetermined, so $\hat u_0=0$. By Lemma~\ref{lem:theta},
$\hat\theta_0=ad_0$. By \eqref{eq:l_theta},
$\hat v_0=\hat\theta_0+\hat u_0=ad_0$.

\emph{Part 2.} By Lemma~\ref{lem:theta}, $\hat\theta_t=a\rho^td_0=-\lvert a\rvert\rho^td_0$.
Substituting into \eqref{eq:l_u} with $d_t=\rho^td_0$,
\[
	\hat u_{t+1}=\lambda\hat u_t+\frac{\bar f}{\bar\tau}\rho^td_0+\omega\lvert a\rvert\rho^td_0
	=\lambda\hat u_t+\Bigl(\frac{\bar f}{\bar\tau}+\omega\lvert a\rvert\Bigr)\rho^td_0 .
\]
With $\hat u_0=0$, iterating gives
\begin{align*}
	\hat u_1&=\Bigl(\tfrac{\bar f}{\bar\tau}+\omega\lvert a\rvert\Bigr)d_0,\\
	\hat u_2&=\Bigl(\tfrac{\bar f}{\bar\tau}+\omega\lvert a\rvert\Bigr)(\lambda+\rho)d_0,
\end{align*}
and in general the expression stated, with $g_t=\sum_{j=0}^{t-1}\lambda^{t-1-j}\rho^j$.
Since $\lambda,\rho\ge0$ and $g_1=1$, $g_t>0$ for $t\ge1$, and the coefficient in
parentheses is positive.

\emph{Part 3.} By \eqref{eq:l_theta}, $\hat v_1=\hat\theta_1+\hat u_1$. With
$\hat\theta_1=a\rho d_0=-\lvert a\rvert\rho d_0$ and $\hat u_1$ from Part 2,
\[
	\hat v_1=-\lvert a\rvert\rho d_0+\Bigl(\frac{\bar f}{\bar\tau}+\omega\lvert a\rvert\Bigr)d_0
	=\Bigl[\frac{\bar f}{\bar\tau}-\lvert a\rvert(\rho-\omega)\Bigr]d_0,
\]
which is negative if and only if \eqref{eq:H1} holds.

\emph{Part 4.} By \eqref{eq:l_theta} and Part 2, for $t\ge1$,
\[
	\hat v_t=-\lvert a\rvert\rho^td_0+\Bigl(\frac{\bar f}{\bar\tau}+\omega\lvert a\rvert\Bigr)g_t\,d_0 .
\]
For $\rho>0$, divide by $\rho^td_0>0$. Using $g_t/\rho^t=\rho^{-1}\sum_{i=0}^{t-1}(\lambda/\rho)^i$,
\[
	\hat v_t<0
	\iff
	\lvert a\rvert>\Bigl(\frac{\bar f}{\bar\tau}+\omega\lvert a\rvert\Bigr)\frac{1}{\rho}
	\sum_{i=0}^{t-1}\Bigl(\frac{\lambda}{\rho}\Bigr)^i .
\]
The right side is strictly increasing in $t$. So once the inequality fails, it fails at all
later dates. This is the first claim.
\begin{itemize}
	\item If $\rho\le\lambda$, the sum diverges, so $\hat v_t$ eventually turns positive.
	\item If $\rho>\lambda$, the right side converges from below to
	$(\bar f/\bar\tau+\omega\lvert a\rvert)/(\rho-\lambda)$. Then $\hat v_t<0$ for every $t$
	if and only if $\lvert a\rvert(\rho-\lambda)\ge\bar f/\bar\tau+\omega\lvert a\rvert$,
	that is, $\lvert a\rvert(\rho-\lambda-\omega)\ge\bar f/\bar\tau$.
\end{itemize}
For $\rho=0$, $\hat v_t=\hat u_t>0$ for $t\ge1$.

\emph{Part 5.} Write $\kappa_u\equiv\bar f/\bar\tau+\omega\lvert a\rvert$. Since
$\tilde u_t\approx\bar u\,\hat u_t$ and $\tilde v_t\approx\bar v\,\hat v_t$,
\[
	\tilde u_t=\bar u\,\kappa_u d_0\,g_t,\qquad
	\tilde v_t=\bar v\bigl(a\rho^t+\kappa_u g_t\bigr)d_0,
\]
with $g_0=0$. For $\rho\ne\lambda$, $g_t=(\rho^t-\lambda^t)/(\rho-\lambda)$. The case
$\rho=\lambda$ follows by continuity. Summing geometric series,
\[
	\sum_{t\ge0}\rho^tg_t=\frac{\rho}{(1-\rho^2)(1-\rho\lambda)},\qquad
	\sum_{t\ge0}g_t^2=\frac{1+\rho\lambda}{(1-\rho^2)(1-\lambda^2)(1-\rho\lambda)} .
\]
The first uses $\frac{1}{1-\rho^2}-\frac{1}{1-\rho\lambda}
=\frac{\rho(\rho-\lambda)}{(1-\rho^2)(1-\rho\lambda)}$. The second uses
\[
	(1-\rho\lambda)(1-\lambda^2)-2(1-\rho^2)(1-\lambda^2)+(1-\rho^2)(1-\rho\lambda)
	=(\rho-\lambda)^2(1+\rho\lambda).
\]
Therefore
\[
	\frac{\sum_t\tilde v_t\tilde u_t}{\sum_t\tilde u_t^2}
	=\frac{\bar v\,\bar u\,\kappa_u d_0^2\bigl(a\sum\rho^tg_t+\kappa_u\sum g_t^2\bigr)}
	{\bar u^2\kappa_u^2d_0^2\sum g_t^2}
	=\bar\theta\Bigl[1+\frac{a}{\kappa_u}\frac{\sum\rho^tg_t}{\sum g_t^2}\Bigr]
	=\bar\theta\Bigl[1-\frac{\lvert a\rvert}{\kappa_u}\frac{\rho(1-\lambda^2)}{1+\rho\lambda}\Bigr],
\]
which is \eqref{eq:slope}. It is negative if and only if
$\lvert a\rvert\rho(1-\lambda^2)>\kappa_u(1+\rho\lambda)$. Substituting
$\kappa_u=\bar f/\bar\tau+\omega\lvert a\rvert$ and rearranging gives \eqref{eq:HP}.

\emph{Corollary~\ref{cor:iid}.} With $\rho=0$, $\hat\theta_t=a\cdot0^t=0$ for $t\ge1$, so
$\hat v_t=\hat u_t$ by \eqref{eq:l_theta}. In \eqref{eq:slope} the correction term has the
factor $\rho=0$, so the slope is $\bar\theta$.
\end{proof}

\section{Interpretation}

\paragraph{Why tightness falls}
At $\rho=0$, \eqref{eq:l_jcc} reads
$\eta_L\hat\theta_0=-[1/(r+\bar\delta)+1/(1-\bar\delta)]d_0$.
\begin{itemize}
	\item The vacancy's flow value, and hence the cost of a hire, rises by $1/(r+\bar\delta)$
	in proportion. That is 167 at D1.
	\item The value of a hire falls only through survival, by $1/(1-\bar\delta)\approx1$.
	\item The expected time to fill, $1/q$, must therefore fall sharply. So tightness falls.
\end{itemize}

\paragraph{The iid case}
Tightness and vacancies fall only on impact, while unemployment is predetermined. At $t=1$
the flow value is back to $\bar K$ and $\theta_1=\bar\theta$. Unemployment rises, and
$v_1=\bar\theta u_1$ rises with it. Entry refills the stock. In the $(u,v)$ plane the
economy drops straight down, jumps up and to the right, and returns along the ray
$v=\bar\theta u$.

At D1, of $\hat u_1=0.16\%$ per 1\% shock:
\begin{itemize}
	\item $0.04$ comes from destruction, $(\bar f/\bar\tau)d_0$;
	\item $0.12$ comes from the date-0 collapse in hiring, $\omega\lvert\hat\theta_0\rvert$.
\end{itemize}

\paragraph{The persistent case}
Tightness stays low while destruction stays high, so vacancies stay below steady state
while unemployment builds up. Once tightness has recovered enough, $v=\theta u$ turns
positive and stays positive.

\paragraph{Where $\bar\delta$ relative to $\bar\tau$ enters}
Conditions \eqref{eq:H1} and \eqref{eq:HP} compare two sensitivities to the same one-point
rise in destruction:
\begin{itemize}
	\item tightness, $\lvert a\rvert$, which is roughly proportional to $1/(r+\bar\delta)$;
	\item unemployment, $\bar f/\bar\tau$.
\end{itemize}
Every other weight is close across calibrations (Table~\ref{tab:weights}). So the
comparison is between $r+\bar\delta$ and $\bar\tau$.
\begin{itemize}
	\item At D1, $(r+\bar\delta)/\bar\tau=0.19$.
	\item With $\bar\delta=\bar\tau$ (CK, $s=0$), $(r+\bar\delta)/\bar\tau=1.11$. The vacancy
	dies with the match, so its flow value cannot be the more sensitive of the two.
\end{itemize}

\section{Numbers}

The calibration uses $\kappa=0$ and $\rho=0.592$, and responses are per 1\% rise in
$\delta_0$, in percent. At D1, $a=-292.0$ (iid: $-280.4$). At $\bar\delta=\bar\tau$,
$a=-47.2$ (iid: $-50.3$).

\begin{table}[h]
\centering
\begin{tabular}{rrrrrrr}
\toprule
 & \multicolumn{3}{c}{D1} & \multicolumn{3}{c}{CK-like ($\bar\delta=\bar\tau$)} \\
\cmidrule(lr){2-4}\cmidrule(lr){5-7}
$t$ & $\hat\theta_t$ & $\hat u_t$ & $\hat v_t$ & $\hat\theta_t$ & $\hat u_t$ & $\hat v_t$ \\
\midrule
0 & $-0.79$ & 0.00 & $-0.79$ & $-1.46$ & 0.00 & $-1.46$ \\
1 & $-0.47$ & 0.17 & $-0.30$ & $-0.87$ & 0.66 & $-0.20$ \\
2 & $-0.28$ & 0.19 & $-0.09$ & $-0.51$ & 0.76 & $0.25$ \\
3 & $-0.16$ & 0.16 & $0.00$ & $-0.30$ & 0.66 & $0.36$ \\
4 & $-0.10$ & 0.13 & $0.03$ & $-0.18$ & 0.51 & $0.33$ \\
\bottomrule
\end{tabular}
\caption{Impulse responses to a persistent destruction shock ($\rho=0.592$, $\kappa=0$).}
\label{tab:irf}
\end{table}

\begin{table}[h]
\centering
\begin{tabular}{lccccc}
\toprule
 & \eqref{eq:H1} & slope, iid & slope, $\rho=0.592$ & $\rho$ needed for slope $<0$ & $(r+\bar\delta)/\bar\tau$ \\
\midrule
D1 & holds & 0.51 & $-0.24$ & $\rho>0.37$ & 0.19 \\
CK-like & holds & 0.51 & $0.17$ & none & 1.11 \\
\bottomrule
\end{tabular}
\caption{Conditions of Proposition~\ref{prop:fe} and the path slope \eqref{eq:slope}.}
\label{tab:slopes}
\end{table}

Condition \eqref{eq:H1} holds in both economies, so the sign at $h=1$ does not separate
them. The path condition \eqref{eq:HP} does.

\section{Remarks}

\begin{remark}\label{rem:flow}
\emph{Holding $K$ fixed.} If the vacancy cost were a constant flow, rather than the flow
value of an asset that $\delta_t$ depreciates, the $1/(r+\bar\delta)$ terms would drop out
of \eqref{eq:l_jcc} and \eqref{eq:a}. Only the survival effect would remain, and
$\lvert a\rvert=1/\{(1-\bar\delta)(\eta_L-B\rho[(1-\bar s)\eta_L-\phi\bar f])\}$ is small.
At our calibration the path slope is then positive in every economy (about $0.49$ at
$\rho=0.592$). This is the flow-cost DMP case.
\end{remark}

\begin{remark}\label{rem:kappa}
With a per-match cost $\kappa>0$, the flow-value terms in \eqref{eq:l_jcc} are scaled by
$x_v=(\bar K/\bar q)/(\kappa+\bar K/\bar q)$. At $x_v=0.5$ the path slopes at $\rho=0.592$
are $-0.23$ (D1) and $0.17$ ($\bar\delta=\bar\tau$).
\end{remark}

\begin{remark}
A permanent shock ($\rho\to1$) recovers the steady-state comparison. It hides
Corollary~\ref{cor:iid}: without persistence, the response is positively sloped after
impact.
\end{remark}

\begin{remark}
\texttt{prop:ds\_asymmetry} Part~3 in the draft gives ``$\bar\delta_e/\bar\tau$ small'' as a
sufficient condition. In the iid free-entry case, entry refills the stock at $t+1$, and
that condition fails for every $\bar\delta_e/\bar\tau$.
\end{remark}

\end{document}
